% Ampliación aditiva: acción, área, memoria sectorial y termometría.
% Propietarios: 70_horizontes_e_informacion.tex; sección térmica antecedente;
% 71_energia_libre_holografica_cerrada.tex del corpus integral.
% Se mantienen las condiciones de la realización cosmológica declarada.

\subsection{Conversion thermique entre le calendrier et l'horizon}
\label{viii:subsec:conversion-termica-calendario-horizonte}

L'action, l'incrément d'aire et la mémoire angulaire admettent une composition énergétique explicite. On conserve une même orientation d'action, la branche positive de Barbero--Immirzi et la réalisation de \eqref{vii:eq:horizonte-cosmologico}. Nous abrégeons
\[
 L=\ell_P=L_\alpha,\qquad
 E_P=\frac{h}{108t_0}=\frac{\hbar c}{L},\qquad
 q_{\mathcal A}=\delta\mathcal A=\gamma_{\rm BI}a_0L^2.
\]
Le symbole \(q_{\mathcal A}\) désigne ici un incrément d'aire.
Le facteur angulaire \(a_0\), la fonctionnelle \(\gamma_{\rm BI}\) et
l'action sont les grandeurs obtenues dans les sections précédentes.
Leur domaine conserve l'incidence hexadique--octadique, l'orientation et
le calendrier du même état enrichi APP--TRIT--TPK.

\begin{proposition}[Composition des échelles surfacique et thermique]
\label{viii:prop:composicion-area-termica}
Avec la tension de~\eqref{vii:eq:tension-informacion-area}, on a
\begin{equation}
 \frac{\mathcal T_{I/A}q_{\mathcal A}}{\log3}
 =\frac{E_P}{\log3}
 =\frac{h}{108t_0\log3}
 =k_B\Theta_{\rm clk}.
 \label{viii:eq:composicion-area-termica}
\end{equation}
Pour \(R=\zeta L\), avec \(\zeta>0\), soit
\(\tau_R=k_BT_{\rm cos}(R)\). Alors
\begin{equation}
 \tau_R=\frac{E_P}{2\pi\zeta},\qquad
 \frac{T_{\rm cos}(R)}{\Theta_{\rm clk}}
 =\frac{\log3}{2\pi\zeta}.
 \label{viii:eq:conversion-dos-temperaturas}
\end{equation}
En particulier, le quotient dans \(R=L\) est \(\log3/(2\pi)\).
\end{proposition}
\begin{proof}
La définition \(\mathcal T_{I/A}=E_P/q_{\mathcal A}\) et la
normalisation tritique de~\eqref{vii:eq:unidad-accion-informacion}
donnent~\eqref{viii:eq:composicion-area-termica}. L'expression
gravitationnelle de la même tension conduit au même résultat :
\[
 \frac{c^4}{\gamma_{\rm BI}a_0GL}
        \gamma_{\rm BI}a_0L^2
 =\frac{c^4L}{G}=\frac{\hbar c}{L}=E_P,
\]
où l'on a utilisé \(G\hbar=c^3L^2\).
D'après~\eqref{vii:eq:horizonte-cosmologico},
\(\tau_R=\hbar c/(2\pi R)=E_P/(2\pi\zeta)\).
Diviser cette égalité par \(k_B\Theta_{\rm clk}=E_P/\log3\)
prouve le quotient des températures.
\end{proof}

Dans la même section d'action, la composition gravitationnelle s'exprime par
\begin{equation}
 Gk_B\Theta_{\rm clk}\log3=GE_P=c^4L,
 \qquad G\tau_R=\frac{c^4L}{2\pi\zeta}.
 \label{viii:eq:composicion-termica-gravitatoria}
\end{equation}
Les deux identités découlent de \(G\hbar=c^3L^2\) et des échelles précédentes.

La tension convertit l'incrément surfacique en énergie de décision.
La normalisation tritique et la périodicité cosmologique déterminent deux
sections thermiques de cette même énergie, reliées par
\eqref{viii:eq:conversion-dos-temperaturas}. La conversion conserve
les deux normalisations et la fonction de Barbero--Immirzi dans la résolution
géométrique.

\subsection{Réponse énergétique de la mémoire sectorielle}
\label{viii:subsec:respuesta-energetica-sectorial}

Fixons \(R>0\) et l'état fidèle
\(\rho_0=\rho_\partial(\Delta_0)\), avec \(\Delta_0>0\),
de~\eqref{vii:eq:estado-sectorial}. Son opérateur modulaire est
\(\mathsf K_0=(\log Z_0)I+\Delta_0\mathsf D_\partial\).
La réalisation d'horizon associe à chaque état de densité
\(\sigma\) les grandeurs de réponse
\begin{equation}
\begin{aligned}
 \mathcal E_R(\sigma)
 &=E_\Lambda(R)+\tau_R
   \operatorname{Tr}[(\sigma-\rho_0)\mathsf K_0],\\
 \mathcal S_R(\sigma)
 &=S_{\rm H}(R)+k_B[s(\sigma)-s(\rho_0)].
\end{aligned}
 \label{viii:eq:respuesta-horizonte-definida}
\end{equation}
Il s'agit de la réponse de l'état modulaire dans la réalisation géométrique déjà spécifiée. Le rayon et l'action fixent son échelle énergétique \(\tau_R\) ; les occupations conservent la provenance de chaque canal.

\Needspace{28\baselineskip}
\begin{theorem}[Hamiltonien de réponse et bilan informationnel]
\label{viii:thm:hamiltoniano-respuesta-termica}
L'opérateur des différences énergétiques correspondant à \eqref{viii:eq:respuesta-horizonte-definida} peut être choisi comme
\begin{equation}
 \mathsf H_R=\tau_R\Delta_0\mathsf D_\partial.
 \label{viii:eq:hamiltoniano-respuesta}
\end{equation}
Son état de Gibbs à l'échelle énergétique \(\tau_R\) est \(\rho_0\).
Pour tout état de densité \(\sigma\) dans l'espace sectoriel,
\begin{equation}
 \mathcal E_R(\sigma)-T_{\rm cos}(R)\mathcal S_R(\sigma)
 =\tau_R\mathcal D(\sigma\Vert\rho_0)\geq0.
 \label{viii:eq:balance-libre-termico}
\end{equation}
Dans une variation différentiable des états autour de \(\rho_0\),
avec \(R\) et \(\mathsf H_R\) fixés, on obtient
\begin{equation}
 \delta\mathcal E_R
 =\tau_R\,\delta s
 =T_{\rm cos}(R)\,\delta\mathcal S_R.
 \label{viii:eq:primera-ley-energetica}
\end{equation}
\end{theorem}
\begin{proof}
Les traces de \(\sigma\) et de \(\rho_0\) valent un, de sorte que la différence élimine le terme scalaire \(\log Z_0\). Il en résulte \(\mathcal E_R(\sigma)-E_\Lambda(R)
=\operatorname{Tr}[(\sigma-\rho_0)\mathsf H_R]\). En outre,
\[
 \frac{e^{-\mathsf H_R/\tau_R}}
 {\operatorname{Tr}e^{-\mathsf H_R/\tau_R}}
 =\frac{e^{-\Delta_0\mathsf D_\partial}}{Z_0}=\rho_0.
\]
L'égalité \(E_\Lambda=T_{\rm cos}S_{\rm H}\) annule la
référence géométrique dans le bilan. L'identité finie
\eqref{vii:eq:ley-modular-finita}, écrite sous la forme
\(s(\sigma)-s(\rho_0)
=\operatorname{Tr}[(\sigma-\rho_0)\mathsf K_0]
-\mathcal D(\sigma\Vert\rho_0)\), démontre
\eqref{viii:eq:balance-libre-termico}. Son terme tangent est la
première loi modulaire déjà prouvée et produit
\eqref{viii:eq:primera-ley-energetica}.
\end{proof}

Pour un incrément unitaire d'occupation, les sauts énergétiques sont
\(\epsilon_{90}=90\tau_R\Delta_0\) et
\(\epsilon_{120}=120\tau_R\Delta_0\). Dans les deux secteurs,
\begin{equation}
 \frac{\epsilon_m}{m\Delta_0}=\tau_R,
 \qquad m\in\{90,120\}.
 \label{viii:eq:pendiente-termica-comun}
\end{equation}
Ainsi, le quotient \(4/3\) entre poids surfaciques enregistre les différents sauts du même hamiltonien. Le quotient entre un saut énergétique et la variation correspondante de \(-\log p\) conserve une échelle commune. La normalisation de l'état n'élimine qu'une constante scalaire.

\begin{corollary}[Susceptibilité énergétique à hamiltonien fixé]
\label{viii:cor:respuesta-escala-termica}
Pour \(\rho_\tau=e^{-\mathsf H_R/\tau}/
\operatorname{Tr}e^{-\mathsf H_R/\tau}\), avec \(\tau>0\),
\begin{equation}
 \frac{\mathrm d\langle\mathsf H_R\rangle_\tau}{\mathrm d\tau}
 =\frac{\operatorname{Var}_{\rho_\tau}(\mathsf H_R)}{\tau^2}>0.
 \label{viii:eq:susceptibilidad-termica}
\end{equation}
En \(\tau=\tau_R\), le résultat est
\(\Delta_0^2\operatorname{Var}_{\rho_0}(\mathsf D_\partial)\).
\end{corollary}
\begin{proof}
La dérivation de la somme finie d'exponentielles donne la variance divisée
par \(\tau^2\). L'état est fidèle et \(\mathsf H_R\) possède plus d'
une valeur propre, de sorte que la variance est positive. La substitution de
\eqref{viii:eq:hamiltoniano-respuesta} donne l'évaluation indiquée.
\end{proof}

Cette dérivée fait varier l'échelle de Gibbs et maintient fixe l'hamiltonien.
La variation radiale de~\eqref{vii:eq:diferencial-horizonte} modifie
également \(\mathsf H_R\) et conserve ses propres termes géométriques.

\subsection{Section thermométrique et conservation du registre}
\label{viii:subsec:seccion-termometrica-registro}

L'état \(\rho_0\) et l'hamiltonien de réponse permettent de fixer une
coordonnée thermométrique antérieure à tout lecteur numérique candidat.
Pour le rayon fixé, soit \(\mathcal L_\Theta=\mathbb R u_R\) une
droite orientée dont la section positive \(u_R\) correspond à l'état
de référence. Dans les états fidèles de Gibbs du même hamiltonien,
\begin{equation}
 \rho_\beta=\frac{e^{-\beta\mathsf H_R}}
                  {\operatorname{Tr}e^{-\beta\mathsf H_R}},\qquad
 \Theta_R(\rho_\beta)=\frac{1}{\beta\tau_R}u_R,
 \quad \beta>0.
 \label{viii:eq:lector-termometrico}
\end{equation}
L'identité opératorielle \(-\log\rho_\beta=\beta\mathsf H_R+(\log Z_\beta)I\) détermine \(\beta\) : soustraire deux valeurs propres d'énergies différentes donne leur quotient unique. En particulier, \(\Theta_R(\rho_0)=u_R\). La droite \(\mathcal L_E\) est l'espace unidimensionnel des grandeurs énergétiques.

\Needspace{23\baselineskip}
\begin{theorem}[Transducteur positif et critère d'identification]
\label{viii:thm:identificacion-transductor}
Il existe une unique application linéaire positive
\(B_R:\mathcal L_\Theta\to\mathcal L_E\) avec
\(B_R(u_R)=\tau_R\). Elle satisfait
\begin{equation}
 B_R(\Theta_R(\rho_\beta))=\beta^{-1}.
 \label{viii:eq:transductor-gibbs}
\end{equation}
La section chronologique compatible est
\begin{equation}
 \Theta_{{\rm clk},R}
 =\frac{E_P}{\tau_R\log3}u_R
 =\frac{2\pi R}{L\log3}u_R,
 \qquad
 B_R(\Theta_{{\rm clk},R})=
 \frac{\mathcal T_{I/A}q_{\mathcal A}}{\log3}.
 \label{viii:eq:seccion-cronologica-compatible}
\end{equation}
Pour toute autre application linéaire \(B_*\) entre les mêmes droites,
\begin{equation}
 B_*=B_R
 \quad\Longleftrightarrow\quad
 B_*(\Theta_{{\rm clk},R})=\frac{h}{108t_0\log3}.
 \label{viii:eq:criterio-una-seccion}
\end{equation}
\end{theorem}
\begin{proof}
La section \(u_R\) est une base d'une droite, de sorte que prescrire
son image définit une unique application linéaire ; la positivité de \(\tau_R\)
garantit la positivité de l'application. Substituer
\eqref{viii:eq:lector-termometrico} prouve
\eqref{viii:eq:transductor-gibbs}. L'expression de \(\tau_R\) et
\eqref{viii:eq:composicion-area-termica} donnent
\eqref{viii:eq:seccion-cronologica-compatible}. Enfin,
\(\Theta_{{\rm clk},R}\) est non nulle. Deux applications linéaires d'une
droite coïncident si et seulement si elles coïncident sur cette section.
\end{proof}

Le critère compare des lecteurs dans une section obtenue préalablement à partir de l'état, de la réponse énergétique et de la normalisation tritique. L'application à une coordonnée numérique concrète exige d'évaluer ce lecteur dans la section spécifiée. Les bases dimensionnelles sont transportées conjointement avec les applications et avec les coordonnées thermiques qu'elles expriment.

\begin{proposition}[Transport de la réponse avec mémoire]
\label{viii:prop:transporte-termico-memoria}
Soit \(J\) une isométrie qui conserve les étiquettes complètes du
registre, et considérons son image comme espace récepteur. Pour
\(\mathsf H'_R=J\mathsf H_RJ^*\), \(\rho'_0=J\rho_0J^*\)
et \(\sigma'=J\sigma J^*\), les espérances énergétiques,
les entropies et le bilan~\eqref{viii:eq:balance-libre-termico} sont conservés.
On a également \(J\mathsf H_R=\mathsf H'_RJ\).
\end{proposition}
\begin{proof}
Sur l'image, \(J\) est unitaire et transporte le calcul fonctionnel.
La cyclicité de la trace conserve les espérances ; les mêmes valeurs propres
non nulles conservent l'entropie et l'entropie relative. L'identité
d'entrelacement utilise \(J^*J=I\). Substituer ces égalités dans le bilan
prouve sa conservation.
\end{proof}

La même équivalence admet une autre échelle \(\tau'=a\tau_R\), avec \(a>0\), et une autre origine énergétique \(e'\). Si \(\mathsf H'=e'I+aJ\mathsf H_RJ^*\), alors
\[
 \frac{\mathsf H'-e'I}{\tau'}
 =J\frac{\mathsf H_R}{\tau_R}J^*,\qquad
 \frac{e^{-\mathsf H'/\tau'}}
      {\operatorname{Tr}_{\operatorname{im}J}e^{-\mathsf H'/\tau'}}
 =J\rho_0J^*.
\]
Le facteur scalaire \(e^{-e'/\tau'}\) s'annule dans l'état ;
les différences énergétiques sont multipliées par \(a\), et l'entropie
relative reste invariante. Le bilan correspondant a donc
pour coefficient \(\tau'\). L'espérance énergétique complète
conserve l'origine : \(\langle\mathsf H'\rangle
=e'+a\langle\mathsf H_R\rangle\). La comparaison des réalisations
s'exprime ainsi au moyen de leurs exposants centrés et de leurs échelles propres.
Si l'espérance alimente les lectures \(m=E/c^2\) et
\(r_g=GE/c^4\), maintenir \(c,G\) donne
\[
 m'=\frac{e'}{c^2}+am,\qquad
 r'_g=\frac{Ge'}{c^4}+ar_g.
\]
La normalisation statistique et le transport des grandeurs maintiennent donc des fonctions distinctes pour la même origine enregistrée.

La composition relie l'incrément surfacique de Barbero--Immirzi, l'énergie du calendrier et la distribution angulaire du bord. Son transport retient la provenance sectorielle et l'archive : le retour de phase reste accompagné de l'avancée de la mémoire de l'état enrichi.
